\subsection{向量空间的子空间的维数与余维数}

命题4. 向量空间E的每个子空间F都是E的一个直因子，并且

\[
\dim \mathbf {F} + \dim (\mathbf {E} / \mathbf {F}) = \dim \mathbf {E}.\tag{3}
\]

由于商向量空间 E/F 是自由模，我们知道（§ 1，第 11 号，命题 21）F 是 E 的一个直因子；于是关系式 (3) 是公式 (1)（第 2 号）的一个特例。

推论 1. 若 E、F、G 是域 K 上的向量空间，则每一个线性映射的正合序列 \(0 \rightarrow E \rightarrow F \rightarrow G \rightarrow 0\) 分裂。

这是表达命题4（§1，第9号）的另一种方式。

推论2. 设 \((\mathbf{E}_i)_{0\leqslant i\leqslant n}\) 是域 K 上向量空间的一个有限族。若存在线性映射的正合序列

\[
0 \longrightarrow \mathrm{E} _ {0} \stackrel {u _ {0}} {\longrightarrow} \mathrm{E} _ {1} \stackrel {u _ {1}} {\longrightarrow} \mathrm{E} _ {2} \longrightarrow \dots \longrightarrow \mathrm{E} _ {n - 1} \stackrel {u _ {n - 1}} {\longrightarrow} \mathrm{E} _ {n} \stackrel {u _ {n}} {\longrightarrow} 0\tag{4}
\]

则关系

\[
\sum_ {2 k + 1 \leqslant n} \dim \mathrm{E} _ {2 k + 1} = \sum_ {2 k \leqslant n} \dim \mathrm{E} _ {2 k}\tag{5}
\]

成立，或者，若所有空间都是有限维的，

\[
\sum_ {i = 1} ^ {n} (- 1) ^ {i} \dim \mathrm{E} _ {i} = 0.\tag{6}
\]

设 \(\mathbf{I}_k = \operatorname{Im} u_k = \operatorname{Ker} u_{k+1}\) 对于 \(0 \leqslant k \leqslant n-1\) ; \(\mathbf{I}_{k+1}\) 因此同构于 \(\mathbf{E}_{k+1}/\mathbf{I}_k\) ，从而（公式（3）） \(\dim \mathbf{I}_k + \dim \mathbf{I}_{k+1} = \dim \mathbf{E}_{k+1}\) 对于 \(0 \leqslant k \leqslant n-2\) 并且此外 \(\dim \mathbf{I}_0 = \dim \mathbf{E}_0\) 和 \(\mathbf{I}_{n-1} = \mathbf{E}_n\) ，因此 \(\dim \mathbf{I}_{n-1} = \dim \mathbf{E}_n\) 。将 \(\dim \mathbf{E}_i\) 用它作为下列和式

\[
\sum_ {k = 0} ^ {n - 1}
\]

dim \(I_{k}\) 的函数所成的表达式代入（5）的两边，我们在每边都得到 \(\sum_{k=0}^{n-1}\dim I_{k}\) ，由此得出推论。

推论3. 若M和N是向量空间E的两个子空间，则

\[
\dim (\mathbf {M} + \mathbf {N}) + \dim (\mathbf {M} \cap \mathbf {N}) = \dim \mathbf {M} + \dim \mathbf {N}.\tag{7}
\]

只需将推论2应用于正合序列即可。

\[
0 \rightarrow \mathbf {M} \cap \mathbf {N} \rightarrow \mathbf {M} \oplus \mathbf {N} \rightarrow \mathbf {M} + \mathbf {N} \rightarrow 0
\]

（§1，第8号，命题10）考虑到以下事实：

\[
\dim (\mathbf {M} \oplus \mathbf {N}) = \dim \mathbf {M} + \dim \mathbf {N}
\]

(第2号，命题2)。

推论4. 设 \(\mathbf{F}\) 是向量空间 \(\mathbf{E}\) 的一个子空间；则 \(\dim \mathbf{F} \leqslant \dim \mathbf{E}\) ；若 \(\mathbf{E}\) 是有限维的，则关系 \(\dim \mathbf{F} = \dim \mathbf{E}\) 等价于 \(\mathbf{F} = \mathbf{E}\) .

第一个断言由(3)显然成立；进一步，若dim E有限，则关系dim F = dim E蕴含dim(E/F) = 0（由(3)），而维数为0的向量空间约化为0。

推论5. 若向量空间 \(\mathbf{E}\) 是向量子空间的一个族 \((\mathbf{F}_{\iota})\) 之和，则

\[
\dim \mathrm{E} \leqslant \sum_ {\iota} \dim \mathrm{F} _ {\iota}.\tag{8}
\]

若进一步 dim E 有限，则(8)式两边相等当且仅当 E 是族 \((\mathbf{F}_{\iota})\) 的直和.

不等式(8)由(3)以及 E 同构于 \(\bigoplus_{\iota} F_{\iota}\) 的一个商这一事实得出（§ 1，第 7 号，公式 (28)）。第二个断言是 § 1，第 10 号，命题 16 的推论 5 的一个特例，因为 (8) 两边的相等意味着除有限多个指标外 \(\dim F_{\iota} = 0\) 。

定义2. 给定向量空间E，E的子空间F的余维数（相对于E），记为 \(\operatorname{codim}_{E} F\) ，或简记为 codim F，即 E/F 的维数（等于 F 在 E 中任一补空间的维数）。

于是关系式 (3) 可写成

\[
\dim \mathrm{F} + \operatorname{codim} \mathrm{F} = \dim \mathrm{E}.\tag{9}
\]

命题5. 设 \(\mathbf{F}\) , \(\mathbf{F}'\) 是向量空间 \(\mathbf{E}\) 的两个子空间，且满足 \(\mathbf{F} \subset \mathbf{F}'\) 。则 \(\operatorname{codim}_{\mathbf{E}} \mathbf{F}' \leqslant \operatorname{codim}_{\mathbf{E}} \mathbf{F} \leqslant \dim \mathbf{E}\) 。若 \(\operatorname{codim}_{\mathbf{E}} \mathbf{F}\) 是有限的，关系式 \(\operatorname{codim}_{\mathbf{E}} \mathbf{F}' = \operatorname{codim}_{\mathbf{E}} \mathbf{F}\) 蕴含 \(\mathbf{F} = \mathbf{F}'\) .

不等式 \(\operatorname{codim}_{\mathbf{E}} \mathbf{F} \leqslant \dim \mathbf{E}\) 由(9)式显然可得；且若 \(\dim \mathbf{E}\) 有限，则关系式 \(\operatorname{codim}_{E} F = \dim E\) 蕴含 \(\dim F = 0\) ，从而 \(F = \{0\}\) 。命题的其余部分由此得出，因为

\[
\operatorname{codim} _ {\mathbf {E}} \mathbf {F} ^ {\prime} = \operatorname{codim} _ {\mathbf {E} / \mathbf {F}} (\mathbf {F} ^ {\prime} / \mathbf {F}),
\]

因为 \(\mathbf{E} / \mathbf{F}'\) 典范同构于 \((\mathbf{E} / \mathbf{F}) / (\mathbf{F}' / \mathbf{F})\) （I，第4节，第7号，定理4）。

命题6. 若M和N是向量空间E的两个子空间，则

\[
\operatorname{codim} (\mathbf {M} + \mathbf {N}) + \operatorname{codim} (\mathbf {M} \cap \mathbf {N}) = \operatorname{codim} \mathbf {M} + \operatorname{codim} \mathbf {N}.\tag{10}
\]

只需将命题4的推论2应用于正合序列即可。

\[
0 \rightarrow \mathrm{E} / (\mathrm{M} \cap \mathrm{N}) \rightarrow (\mathrm{E} / \mathrm{M}) \oplus (\mathrm{E} / \mathrm{N}) \rightarrow \mathrm{E} / (\mathrm{M} + \mathrm{N}) \rightarrow 0
\]

(§ 1，第 7 号，命题 10)并使用第 2 号，命题 2。

注意，如果 \(\mathbf{E}\) 是有限维的，(10) 是 (7) 和 (9) 的推论。

命题7. 若 \((\mathbf{F}_i)\) 是向量空间 \(\mathbf{E}\) 的子空间的一个有限族，则 \(\operatorname{codim}\left(\bigcap_i \mathbf{F}_i\right) \leqslant \sum_i \operatorname{codim} \mathbf{F}_i\) .

若 \(\mathbf{F} = \bigcap_{i}\mathbf{F}_{i}\) ，则 \(\mathbf{E} / \mathbf{F}\) 同构于诸 \(\mathbf{E} / \mathbf{F}_i\) 的直和的一个子空间（§1，第7号，公式（27））。

向量空间 E 的维数为 1（相应地，维数为 2）的向量子空间通常称为过 0 的直线（相应地，过 0 的平面）（若不会引起混淆，则简称为直线（相应地，平面）（参见第 9 节第 3 号）），这是类比经典几何的语言；E 的子空间若其余维数为 1，则称为过 0 的超平面（或简称为超平面）。超平面也可以定义为集合 S 的极大元，其中 S 是按包含关系排序的、E 中不同于 E 的向量子空间全体。包含子空间 H 的 E 的子空间与 E/H 的子空间之间存在一一对应（I，第 4 节，第 7 号，定理 4）；若 E 的维数 ≥1，则 S 非空，且说 H 是 S 中的极大元意味着 E/H 不包含不同于 \{0\} 和 E/H 的子空间，这蕴含 E/H 由它的任意非零元生成，换言之，它的维数为 1。

在维数为 \(n \geqslant 1\) 的有限维向量空间中，由公式 (3)，超平面就是维数为 n - 1 的那些子空间。

命题 8. 在域 K 上的向量空间 E 中，每个向量子空间 F 都是包含它的那些超平面的交集。

只需证明对任意 \(x \notin F\) ，存在一个包含 F 但不包含 x 的超平面 H。由假设 \(F \cap Kx = \{0\}\) ，故 F 与 Kx 的和 M 是直和。设 N 是 M 在 E 中的补子空间；则 E 是 \(H = F + N\) 与 Kx 的直和，因此 H 就是具有所需性质的超平面。

注。本节中对于向量空间的子空间所证明的大部分性质，对于维数（§ 2，第 7 号，注 2）已定义的自由模的子模并不成立。*例如，交换环的理想不一定有基，因为存在某些整环 A，其中某些理想不是主理想（VII，§ 1，第 1 号），且此类环中任意两个元素都是线性相关的（§ 1，第 11 号，注 1）。*自由 A-模 E 的子模可以是自由的，不同于 E，且与 E 具有相同的维数，正如整环 A 中的主理想所示；同样的例子还表明，自由 A-模的自由子模不一定有补子模。

